\naslov{Togo telo}

\begin{definicija}
  Gibanje sistema materialnih točk je \pojem{togo}, če gibanje ohranja razdalje
  med točkami.
\end{definicija}

\begin{definicija}
  Telo je \pojem{togo}, če nima koles.
  Alternativno je telo togo, če je edini njegov način gibanja togo gibanje.
\end{definicija}

\begin{izrek}
  Izometrija je afina preslikava.
\end{izrek}

\begin{proof}
  Izberimo $\mb{P}_0$ in točke $\mb{A}, \mb{B}, \mb{C}$, ki niso kolinearne.
  Naj bo $\mb{P}$ poljubna točka.
  Za $\vec{r} = \mb{P} - \mb{P}_0$, $\vec{a} = \mb{A} - \mb{P}_0$, $\vec{b} =
  \mb{B} - \mb{P}_0$ in $\vec{c} = \mb{C} - \mb{P}_0$ velja
  \[
	\vec{r} = \alpha \vec{a} + \beta \vec{b} + \gamma \vec{c}.
  \]
  Izometrija naj slika $\mb{T}$ v $\mb{T}'$, poleg tega naj je
  \[
	\pvec{r}' = \alpha' \pvec{a}' + \beta' \pvec{b}' + \gamma' \pvec{c}'.
  \]
  Če enačbo za $\vec{r}$ skalarno množimo z $\vec{a}$, $\vec{b}$ oziroma
  $\vec{c}$, dobimo enačbe
  \begin{align*}
	\vec{r} \cdot \vec{a} &= \alpha \vec{a} \cdot \vec{a} + \beta \vec{b} \cdot \vec{a} + \gamma \vec{c} \cdot \vec{a} \\
	\vec{r} \cdot \vec{b} &= \alpha \vec{a} \cdot \vec{b} + \beta \vec{b} \cdot \vec{b} + \gamma \vec{c} \cdot \vec{b} \\
	\vec{r} \cdot \vec{c} &= \alpha \vec{a} \cdot \vec{c} + \beta \vec{b} \cdot \vec{c} + \gamma \vec{c} \cdot \vec{c}
  \end{align*}
  To zložimo v matriko $A$, da dobimo
  \[
	\begin{bmatrix}
	  \vec{r} \cdot \vec{a} \\
	  \vec{r} \cdot \vec{b} \\
	  \vec{r} \cdot \vec{c} \\
	\end{bmatrix}
	= A
	\begin{bmatrix}
	  \alpha \\ \beta \\ \gamma
	\end{bmatrix}
  \]
  in podobno za količine s črto.
  Izometrija ohranja tako razdalje kot kote, torej sta levi strani sistemov in
  matriki $A, A'$ enaki.
  Sledi torej $\alpha = \alpha'$, $\beta = \beta'$ in $\gamma = \gamma'$.
  Dobimo
  \[
	\mb{P}' = \mb{P}_0' + \pvec{r}' = \mb{P}_0' + A \vec{r}
	= \mb{P}_0' + A(\mb{P} - \mb{P}_0)
  \]
  in $A \in O(3)$.
\end{proof}

Togo gibanje lahko zapišemo v obliki
\[
  \mb{P}(t) = \mb{P}_0(t) + Q(t) (\mb{P}(t=0) -\mb{P}_0(t=0)).
\]
Predstavimo ga kot relativno gibanje, kjer je RKS položaj telesa, kot ga vidimo
ob času $t=0$, AKS pa položaj telesa ob času $t$.
RKS se giblje skupaj s telesom, zato mu pravimo \pojem{telesni KS}, absolutnemu
pa pravimo \pojem{prostorski KS}.

\vprasanje{Kako opišeš gibanje togega telesa?}

Za telo $B$ definiramo volumensko in masno mero in predpostavimo, da za $B' =
B(t)$ velja $m(B') = m(B)$.
Sledi
\[
  m(B') = \iiint_{B'} \rho' dV'
  = \iiint_{B} \rho'(\mb{P}'(\mb{P})) \abs{\det \frac{\partial
	  \mb{P}'}{\partial \mb{P}}} dV.
\]
Ker pa je $\mb{P}' = \mb{P}_0' + Q(\mb{P} - \mb{P}_0)$, torej je odvod enak $Q$
in njegova determinanta enaka $1$.
Torej
\[
  \iiint_{B} \rho'(\mb{P}'(\mb{P})) dV
  = m(B')
  = m(B)
  = \iiint_B \rho(\mb{P}) dV,
\]
torej (ker to velja tudi na delih telesa) $\rho'(\mb{P}'(\mb{P})) =
\rho(\mb{P})$.
Podobno izpeljemo, da je
\[
  \mb{P}_*' = \inv{m(B)} \iiint_{B'} \mb{P}' - \mb{P}_0' dm + \mb{P}_0'
\]
enak $\mb{P}_*' = \mb{P}_0' + Q(\mb{P}_* - \mb{P}_0)$ za $\mb{P}_*$, definiran s
podobnim integralom.
Podobno kot v diskretnem primeru lahko tudi tu izpeljemo
\[
  \dot{\mb{P}}_*' = \inv{m(B')} \iiint_{B'} \dot{\mb{P}}' dm.
\]

Vrtilno količino izračunamo kot
\begin{align*}
  \vec{l}'(\mb{P}_0')
  &= \iiint_{B'} (\mb{P}' - \mb{P}_0') \times \partial_t (\mb{P}' - \mb{P}_0') dm' \\
  &= \iiint_{B'} \pvec{\zeta}' \times (\pvec{\omega}' \times \pvec{\zeta}') dm' \\
  &= \iiint_{B'} \abs{\pvec{\zeta}'}^2 \pvec{\omega}' - (\pvec{\zeta}' \cdot \pvec{\omega}') . \pvec{\zeta}' dm' \\
  &= \iiint_{B'} \left(\abs{\pvec{\zeta}'}^2 I - \pvec{\zeta}' \otimes \pvec{\zeta}'\right) dm' \cdot \pvec{\omega}'.
\end{align*}
Iz tega dobimo definicijo prostorskega vztrajnostnega tenzorja
\[
  J'(\mb{P}_0') = \iiint_{B'} \left(\abs{\pvec{\zeta}'}^2 I - \pvec{\zeta}'
	\otimes \pvec{\zeta}'\right) dm'
  = \iiint_{B'} \abs{\mb{P}' - \mb{P}_0'}^2 I - (\mb{P}' - \mb{P}_0') \otimes
  (\mb{P}' - \mb{P}_0') dm'.
\]
Podobno definiramo telesni vztrajnostni tenzor, izračunamo lahko $J' = Q J Q^T$.
Vztrajnostni tenzor je simetričen, poleg tega pa ga lahko diagonaliziramo
\[
  J(\mb{P}_0') = \sum_{i=1}^3 J_i \vec{w}_i \otimes \vec{w}_i,
\]
kjer so $J_i$ \pojem{glavne} (lastne) vrednosti.

\begin{trditev}
  Vztrajnostni tenzor je nenegativen.
  Če $B$ ni tanka palica, je pozitivno definiten.
\end{trditev}

\begin{proof}
  Trdimo, da je $\vec{e} \cdot J \vec{e} \ge 0$ za enotski $\vec{e}$.
  Po Cauchy-Schwarzu velja
  \[
	\iiint_B \abs{\vec{\zeta}}^2 - (\vec{e} \cdot \vec{\zeta})^2 dm \ge 0
  \]
  natanko tedaj, ko je $\vec{e}$ vzporeden $\vec{\zeta}$, kar se zgodi le, če je
  $B$ tanka palica.
\end{proof}

\vprasanje{Definiraj vztrajnostni tenzor in dokaži, da je pozitivno definiten.}

\begin{trditev}
  Normala na ravnino zrcalne simetrije telesa je glavna smer vztrajnostnega
  tenzorja s polom na tej ravnini.
\end{trditev}

\begin{proof}
  Naj bo $\vec{e}$ ta normala.
  Trdimo $J(\mb{P}_0) \vec{e} = J \vec{e}$ za neko lastno vrednost $J$.
  To je ekvivalentno temu, da je deviator
  \[
	D \vec{e} = \vec{\zeta} \otimes \vec{\zeta} \vec{e}
  \]
  vzporeden $\vec{e}$.
  Naj bo $\vec{f} \bot \vec{e}$ in izračunajmo
  \[
	\vec{f} \cdot D \vec{e}
	= \iiint_B \vec{f} \cdot \vec{\zeta} . \vec{e} \cdot \vec{\zeta} dm
	= \iiint_B \vec{f} \cdot \pvec{\zeta}' . \vec{e} \cdot \pvec{\zeta}' dm
  \]
  za zrcalno sliko $\pvec{\zeta}'$.
  Izračunamo lahko, da velja $\vec{e} \cdot \pvec{\zeta}' = - \vec{e} \cdot
  \vec{\zeta}$ in $\vec{f} \cdot \pvec{\zeta}' = \vec{f} \cdot \vec{\zeta}$,
  torej velja $\vec{f} \cdot D \vec{e} = 0$.
\end{proof}

\begin{trditev}
  Naj ima $B$ dve ravnini simetrije z normalama $\vec{e}_1$ in $\vec{e}_2$.
  Potem ima $J(\mb{P}_0)$, kjer je $\mb{P}_0$ v preseku obeh ravnin, glavne
  smeri $\vec{e}_1, \vec{e_2}, \vec{e}_1 \times \vec{e}_2$.
\end{trditev}

\vprasanje{Kaj lahko poveš o vztrajnostnemu tenzorju, če ima eno ali dve ravnini
  zrcalne simetrije? Dokaži.}